% META: {"id": "1SPE-SECDEG-CR-000", "chapitre": "1SPE-SECOND-DEGRE", "type_objet": "cours", "sous_type": "ouverture", "titre": "Fonctions polynomes du second degre", "temps_gabarit": 1, "parcours_duree": {"P1": "12h", "P2": "10h", "P3": "8h"}, "prerequis": ["R1", "R2", "R3", "R4", "R5"], "capacites": ["C1", "C2", "C3", "C4", "C5", "C6"], "mode_creation": "ex_nihilo", "sources_inspiration": [], "status": "generated"}

\section*{Contrat du chapitre}

\subsection*{Ce que tu vas savoir faire}

À la fin de ce chapitre, tu seras capable de :

\begin{itemize}
  \item[\textbf{C1}] Reconnaître un polynôme du second degré et passer d'une forme à l'autre (développée, canonique, factorisée).
  \item[\textbf{C2}] Représenter une parabole, identifier son sommet et son axe de symétrie, et dresser le tableau de variations.
  \item[\textbf{C3}] Calculer le discriminant et résoudre une équation du second degré en distinguant les trois cas selon le signe de $\Delta$.
  \item[\textbf{C4}] Factoriser un trinôme et étudier le signe d'une expression du second degré à l'aide d'un tableau de signes.
  \item[\textbf{C5}] Résoudre une inéquation du second degré en ramenant à l'étude du signe d'un trinôme.
  \item[\textbf{C6}] Modéliser une situation concrète par un polynôme du second degré et déterminer un extremum (optimisation).
\end{itemize}

\subsection*{Ce dont tu as besoin avant de commencer}

Ce chapitre s'appuie sur cinq familles de prérequis. Le diagnostic de la page suivante te dira lesquels tu maîtrises déjà.

\begin{itemize}
  \item[\textbf{R1}] \textbf{Identités remarquables} — développer et factoriser $(a+b)^2$, $(a-b)^2$, $(a+b)(a-b)$.
  \item[\textbf{R2}] \textbf{Équations du premier degré} — résoudre $ax + b = 0$, mettre en évidence un facteur commun.
  \item[\textbf{R3}] \textbf{Produit nul} — utiliser la règle « un produit est nul si et seulement si l'un de ses facteurs est nul ».
  \item[\textbf{R4}] \textbf{Fonctions} — calculer une image, chercher un antécédent, lire un tableau de valeurs ; connaître la fonction carré $x \mapsto x^2$.
  \item[\textbf{R5}] \textbf{Tableaux de signes} — dresser le tableau de signes d'une expression affine, d'un produit ou d'un quotient.
\end{itemize}

Si une fiche R te concerne, lis-la avant d'aborder le cours : chaque notion manquante sera un obstacle au moment où tu en auras le plus besoin.

\subsection*{Temps de travail estimé}

\begin{itemize}
  \item Parcours~\parcoursUn{} (Consolidation) : environ \textbf{12~h}.
  \item Parcours~\parcoursDeux{} (Maîtrise) : environ \textbf{10~h}.
  \item Parcours~\parcoursTrois{} (Approfondissement) : environ \textbf{8~h}.
\end{itemize}

\subsection*{Situation d'accroche — La boîte en carton}

On découpe quatre carrés identiques de côté $x$ (en cm) aux coins d'une feuille de carton rectangulaire de dimensions $30$~cm $\times$ $20$~cm, puis on relève les bords pour former une boîte ouverte.

\begin{itemize}
  \item Exprime le volume $V(x)$ de la boîte en fonction de $x$.
  \item Quelles valeurs de $x$ sont admissibles ?
  \item Pour quelle découpe l'aire du fond de la boîte vaut-elle $336$~cm$^2$ ?
\end{itemize}

On obtient $V(x) = x(30 - 2x)(20 - 2x)$, dont le développement vaut
\textbf{exactement} $4x^3 - 100x^2 + 600x$ : c'est un polynôme de
\textbf{degré 3}, et non du second degré. Chercher la découpe qui rend ce volume
maximal demande un outil que ce chapitre ne fournit pas — la dérivation, étudiée
plus tard ; la valeur cherchée n'est d'ailleurs même pas un nombre décimal.

Le second degré n'en est pas absent pour autant. L'aire du fond,
$A(x) = (30-2x)(20-2x) = 4x^2 - 100x + 600$, est bien un trinôme : la troisième
question se ramène à l'équation $4x^2 - 100x + 600 = 336$, que le discriminant
résout. Ses deux racines sont $3$ et $22$, et seule $x = 3$ appartient au domaine
admissible $]0\,;\,10[$. Les outils de ce chapitre — formes du trinôme,
discriminant, tableau de signes, forme canonique — traitent ainsi toute situation
dont le modèle est quadratique, et c'est cette compétence de modélisation (C6) que
la boîte sert ici à installer.

Le Temps~7 (TD fil rouge) reprend la situation entière : la fonction auxiliaire du
second degré avec les outils du chapitre, et le maximum du volume par une
exploration numérique, faute de la dérivation.

% BEGIN-VERIFY
% from sympy import symbols, expand, diff, solve, Rational, simplify, Poly
%
% x = symbols('x', real=True)
%
% # La situation d'accroche : boite obtenue en decoupant quatre carres de cote
% # x aux coins d'un carton 30 x 20, puis en relevant les bords.
% volume = x * (30 - 2*x) * (20 - 2*x)
% assert simplify(expand(volume) - (4*x**3 - 100*x**2 + 600*x)) == 0
%
% # Domaine : il faut x > 0 et les deux dimensions restantes positives.
% assert 30 - 2*x != 0
% domaine = (0, 10)                        # 20 - 2x > 0 impose x < 10
% assert volume.subs(x, 0) == 0 and volume.subs(x, 10) == 0
%
% # Le maximum se trouve en annulant la derivee, et il est bien dans ]0;10[.
% critiques = [r for r in solve(diff(volume, x), x) if r.is_real]
% dans_domaine = [r for r in critiques if domaine[0] < r < domaine[1]]
% assert len(dans_domaine) == 1
% optimal = dans_domaine[0]
% assert Rational(39, 10) < optimal < Rational(41, 10)
%
% # C'est un maximum GLOBAL sur le domaine, verifie par balayage.
% meilleur = volume.subs(x, optimal)
% for k in range(1, 100):
%     essai = Rational(k, 10)
%     assert volume.subs(x, essai) <= meilleur
%
% # La boite est bien ouverte : sa hauteur est x, sa base (30-2x)(20-2x).
% hauteur, base = optimal, (30 - 2*optimal) * (20 - 2*optimal)
% assert simplify(hauteur * base - meilleur) == 0
%
% # Ce que l'accroche affirme desormais, verifie ici meme.
% # 1. Le developpement du volume est EXACT, ce n'est pas une approximation.
% assert expand(volume) == 4*x**3 - 100*x**2 + 600*x
% assert Poly(volume, x).degree() == 3
%
% # 2. La decoupe optimale n'est pas decimale : ce chapitre ne peut pas la donner.
% assert not optimal.is_rational
%
% # 3. L'aire du fond, elle, est bien un trinome du second degre.
% aire = (30 - 2*x) * (20 - 2*x)
% assert expand(aire) == 4*x**2 - 100*x + 600
% assert Poly(aire, x).degree() == 2
%
% # 4. L'equation A(x) = 336 a pour racines 3 et 22 ; seule 3 est admissible.
% racines = sorted(solve(expand(aire) - 336, x))
% assert racines == [3, 22]
% assert [r for r in racines if domaine[0] < r < domaine[1]] == [3]
% assert aire.subs(x, 3) == 336
% END-VERIFY
